% !TEX root = ../../main.tex
\subsection{Merkle 日志完整性验证机制}
\label{subsec:merkle}

\subsubsection*{叶子节点构造}

Merkle Tree 以订单内全部消息的哈希值列表作为叶子节点，叶子节点 $L_i$ 即为第 $i$ 条消息的 \texttt{message\_hash}：
\[
  L_i = h_i = \operatorname{SM3}(\mathit{orderID} \| \mathit{senderPID}_i \| \mathit{role}_i \| \mathit{content}_{i,\mathrm{plain}} \| \mathit{timestamp}_i)
\]
叶子节点按消息写入顺序排列（\texttt{timestamp ASC, id ASC}），顺序的固定性保证任意消息重排均会导致 Merkle Root 改变。

\subsubsection*{树的构建规则}

Merkle Tree 采用迭代方式自底向上构建：

\begin{enumerate}[label=\arabic*., leftmargin=2em]
  \item 以所有叶子节点哈希值的小写十六进制表示作为第 0 层；
  \item 每层相邻两个节点两两拼接后计算 SM3 哈希，得到父节点：
    \[
      \mathrm{parent}(i) = \operatorname{SM3}(L_{2i} \| L_{2i+1})
    \]
  \item 若当前层节点数为奇数，最后一个节点与自身拼接（复制）后参与父节点计算；
  \item 重复上述过程直至仅剩一个节点，该节点即为 Merkle Root $R$；
  \item 若订单内无消息，Root 取固定值 $\operatorname{SM3}(\text{"EMPTY"})$。
\end{enumerate}

\subsubsection*{快照触发时机}

Merkle Root 快照写入 \texttt{audit\_logs} 表的触发时机包括以下几种：

\begin{itemize}[leftmargin=2em]
  \item \textbf{实时快照}：每条消息发送成功并写入数据库后，服务器立即调用 \texttt{create\_merkle\_snapshot} 函数，以当前全部消息哈希重新计算 Root，将 $\langle\mathit{order\_id},\allowbreak\; \text{MERKLE\_ROOT\_UPDATED},\allowbreak\; \mathit{root},\allowbreak\; \mathit{message\_count}\rangle$ 写入 \texttt{audit\_logs}；
  \item \textbf{管理员手动快照}：管理员可通过 \texttt{POST /admin/snapshot/\{order\_id\}} 接口手动触发某订单的 Merkle Root 快照；
  \item \textbf{历史补全快照}：管理员可通过 \texttt{POST /admin/backfill\_snapshots} 接口对所有尚未生成快照的历史订单批量补全 Root 快照。
\end{itemize}

\subsubsection*{完整性验证流程}

管理员执行日志完整性验证时（\texttt{POST /admin/verify/\{order\_id\}}），系统按以下步骤完成验证：

% !TEX root = ../../../main.tex
% Merkle 验证流程图
\begin{figure}[H]
\centering
\begin{tikzpicture}[
  node distance = 0.55cm and 2.4cm,
  proc/.style     = {rectangle, draw, rounded corners=3pt, minimum width=3.4cm, minimum height=0.65cm,
                     font=\small, fill=blue!8, draw=blue!50!black},
  decision/.style = {diamond, draw, aspect=2.2, minimum width=3.4cm, minimum height=0.85cm,
                     font=\small, align=center, fill=orange!15, draw=orange!60!black, inner sep=2pt},
  term_ok/.style  = {rectangle, draw, rounded corners=8pt, minimum width=3.0cm, minimum height=0.65cm,
                     font=\small, fill=green!14, draw=green!50!black},
  term_fail/.style= {rectangle, draw, rounded corners=8pt, minimum width=3.0cm, minimum height=0.65cm,
                     font=\small, fill=red!12, draw=red!50!black},
  arr/.style      = {-Stealth, thick},
  lbl/.style      = {font=\small},
]

\node[proc]     (read)   {读取 messages 表中消息列表};
\node[proc,     below=0.6cm of read]    (decrypt) {SM4-CBC 解密，获取明文及密文摘要};
\node[proc,     below=0.6cm of decrypt] (rehash)  {逐条重算 SM3 摘要作为叶子节点};
\node[proc,     below=0.6cm of rehash]  (rebuild) {重建 Merkle Tree，计算新 Root'};
\node[proc,     below=0.6cm of rebuild] (load)    {从 audit\_logs 加载历史快照 Root};
\node[decision, below=0.8cm of load]    (d_eq)    {Root'\\== Root？};
\node[term_ok,  right=2.4cm of d_eq]   (ok)      {VERIFY\_OK：日志完整};
\node[term_fail,left=2.4cm of d_eq]    (fail)    {VERIFY\_FAIL：存在篡改};
\node[proc,     below=0.9cm of d_eq]   (log)     {写入审计操作日志};

\draw[arr] (read)    -- (decrypt);
\draw[arr] (decrypt) -- (rehash);
\draw[arr] (rehash)  -- (rebuild);
\draw[arr] (rebuild) -- (load);
\draw[arr] (load)    -- (d_eq);
\draw[arr] (d_eq)    -- node[lbl,above]{是} (ok);
\draw[arr] (d_eq)    -- node[lbl,above]{否} (fail);
\draw[arr] (d_eq)    -- node[lbl,left]{} (log);
\draw[arr] (ok.south)   -- ++(0,-0.3) |- (log.east);
\draw[arr] (fail.south) -- ++(0,-0.3) |- (log.west);

\end{tikzpicture}
\caption{Merkle 日志完整性验证流程图}
\label{fig:merkle_verify}
\end{figure}


\begin{enumerate}[label=\arabic*., leftmargin=2em]
  \item 从数据库按序读取订单内所有消息，对 \texttt{content} 字段执行 SM4 解密，还原明文；
  \item 对每条消息按相同公式重新计算消息哈希 $h_i'$，与数据库中存储的 \texttt{message\_hash} 对比；若存在不匹配的条目，记录差异列表；
  \item 以重新计算的哈希序列 $[h_1', h_2', \ldots, h_n']$ 为叶子节点重建 Merkle Tree，得到当前 Root $R'$；
  \item 从 \texttt{audit\_logs} 读取最近一次 \texttt{MERKLE\_ROOT\_UPDATED} 快照，取其 \texttt{merkle\_root} 字段作为历史快照值 $R_{\mathrm{snap}}$；
  \item 比较 $R' \stackrel{?}{=} R_{\mathrm{snap}}$：若一致且无哈希不匹配条目，则判定完整性验证通过（\texttt{VERIFY\_OK}）；否则判定为验证失败（\texttt{VERIFY\_FAIL}）；
  \item 将验证结果（包括不匹配哈希列表、当前 Root、快照 Root 和消息数量）写入 \texttt{audit\_logs}。
\end{enumerate}

该机制能够检测以下篡改行为：
\begin{itemize}[leftmargin=2em]
  \item 修改某条消息的内容或时间戳：重新计算的 $h_i'$ 将与存储的 \texttt{message\_hash} 不一致，且 $R' \neq R_{\mathrm{snap}}$；
  \item 删除某条消息：叶子节点数量减少，树结构改变，$R' \neq R_{\mathrm{snap}}$；
  \item 插入伪造消息：新增叶子节点，$R' \neq R_{\mathrm{snap}}$；
  \item 重排消息顺序：叶子节点顺序改变，$R' \neq R_{\mathrm{snap}}$。
\end{itemize}
